\subsection{有限次扩张中的范数与迹}

代数中的传递公式（III，第 548 页）在有限次扩张的情形下蕴含下述命题。

命题 2.——设 F 是 K 的有限次扩张，E 是 F 的一个子扩张。则对每个 \(x \in F\) 我们有

\[
\operatorname{Tr} _ {\mathrm{F} / \mathrm{K}} (x) = \operatorname{Tr} _ {\mathrm{E} / \mathrm{K}} (\operatorname{Tr} _ {\mathrm{F} / \mathrm{E}} (x))\tag{6}
\]

\[
\mathrm{N} _ {\mathrm{F} / \mathrm{K}} (x) = \mathrm{N} _ {\mathrm{E} / \mathrm{K}} (\mathrm{N} _ {\mathrm{F} / \mathrm{E}} (x))\tag{7}
\]

推论.——设 \(m = [F : E]\) ；则对每个 \(x \in E\) 我们有

\[
\operatorname{Tr} _ {\mathrm{F} / \mathrm{K}} (x) = m. \operatorname{Tr} _ {\mathrm{E} / \mathrm{K}} (x)\tag{8}
\]

\[
\mathrm{N} _ {\mathrm{F} / \mathrm{K}} (x) = \mathrm{N} _ {\mathrm{E} / \mathrm{K}} (x) ^ {m}\tag{9}
\]

命题 3.——设 E 是 K 的 n 次有限扩张，x 是 E 中在 K 上次数为 d 的元素。记 \(f(X) = X^d + \sum_{i=1}^{d} a_i X^{d-i}\) 为 x 在 K 上的极小多项式。则我们有

\[
\mathrm{Tr} _ {\mathrm{E/K}} (x) = - \frac {n}{d} a _ {1}\tag{10}
\]

\[
\mathrm{N} _ {\mathrm{E} / \mathrm{K}} (x) = ((- 1) ^ {d} a _ {d}) ^ {n / d} = (- 1) ^ {n} a _ {d} ^ {n / d}\tag{11}
\]

命题 3 直接由命题 2 的推论及下述引理得出：

引理 3.——设 R 是交换环，\(f(\mathbf{X}) = \mathbf{X}^d + \sum_{i=1}^{d} a_i \mathbf{X}^{d-i}\) 是 R{[}X{]} 中的一个首一多项式，A 是 R-代数 R{[}X{]}/(f)，x 是 X 在 A 中的剩余类。则 \(\operatorname{Tr}_{\mathrm{A/R}}(x) = -a_1\) 且 \(\mathrm{N}_{\mathrm{A/R}}(x) = (-1)^d a_d\) .

由该推论（IV，第 11 页），序列 \((1, x, \ldots, x^{d-1})\) 是 \(\mathbf{A}\) 的一组基；进一步我们有

\[
\begin{array}{r l} x \cdot 1 = x, & x \cdot x = x ^ {2}, \dots , x \cdot x ^ {d - 2} = x ^ {d - 1}, \\ & x \cdot x ^ {d - 1} = - a _ {d} \cdot 1 - a _ {d - 1} \cdot x - \dots - a _ {1} \cdot x ^ {d - 1}. \end{array}
\]

于是，相对于 A 的基 \((1, x, \ldots, x^{d-1})\) 表示乘以 x 的矩阵具有下述形式（为确定起见取 d = 5）：

\[
\left( \begin{array}{c c c c c} 0 & 0 & 0 & 0 & - a _ {5} \\ 1 & 0 & 0 & 0 & - a _ {4} \\ 0 & 1 & 0 & 0 & - a _ {3} \\ 0 & 0 & 1 & 0 & - a _ {2} \\ 0 & 0 & 0 & 1 & - a _ {1} \end{array} \right).
\]

这个矩阵的迹显然是 \(-a_{1}\) ；把行列式按第一行展开即可算得，于是我们得到

\[
(- 1) ^ {d - 1} (- a _ {d}) = (- 1) ^ {d} a _ {d}.
\]

在本小节其余部分，我们用 E 表示 K 的一个有限次扩张，并用 x 表示 E 中的一个元素。我们将指出在各种情形下如何计算 x 的范数与迹。

\begin{enumerate}
\def\labelenumi{\alph{enumi})}
\tightlist
\item
  可分扩张的情形：设 E 在 K 上是 n 次可分扩张，用 \(\Omega\) 表示 K 的一个代数闭包，用 \(\sigma_{1}, \ldots, \sigma_{n}\) 表示 E 到 \(\Omega\) 中的 n 个互不相同的 K-同态。由公式 (3)（V，第 48 页），我们在 \(\Omega\) 中有
\end{enumerate}

\[
\mathrm{Tr} _ {\mathrm{E} / \mathrm{K}} (x) = \sum_ {i = 1} ^ {n} \sigma_ {i} (x), \quad \mathrm{N} _ {\mathrm{E} / \mathrm{K}} (x) = \prod_ {i = 1} ^ {n} \sigma_ {i} (x).\tag{12}
\]

\begin{enumerate}
\def\labelenumi{\alph{enumi})}
\setcounter{enumi}{1}
\tightlist
\item
  p-根式扩张的情形：设 K 的特征为 p \textgreater{} 0 且扩张 E 是 p-根式的；存在整数 \(e \geqslant 0\) 使得 \([E : K] = p^{e}\) （V，第 26 页，命题 4）。若 f 是 x 在 K 上的高度，则 x 在 K 上的极小多项式是 \(X^{p^{f}} - x^{p^{f}}\) （V，第 24 页，命题 1）。由命题 3，我们有 \(\mathrm{N}_{\mathrm{E}/\mathrm{K}}(x) = (x^{p^{f}})^{p^{e}/p^{f}}\) ，从而
\end{enumerate}

\[
\mathrm{N} _ {\mathrm{E} / \mathrm{K}} (x) = x ^ {p ^ {e}} = x ^ {[ \mathrm{E}: \mathrm{K} ]}.\tag{13}
\]

对于迹，我们求得 \(\mathrm{Tr}_{\mathrm{E} / \mathrm{K}}(x) = -p^{e - f}a\) ，其中 \(a\) 是 \(\mathbf{X}^{p^f - 1}\) 在多项式 \(\mathbf{X}^{p^f} - x^{p^f}\) 中的系数；换言之，我们有

\[
\operatorname{Tr} _ {\mathrm{E} / \mathrm{K}} (x) = p ^ {e} \cdot x = [ \mathrm{E}: \mathrm{K} ] x = \left\{ \begin{array}{l l} x & \text { if } \quad [ \mathrm{E}: \mathrm{K} ] = 1 \\ 0 & \text { if } \quad [ \mathrm{E}: \mathrm{K} ] > 1. \end{array} \right.\tag{14}
\]

\begin{enumerate}
\def\labelenumi{\alph{enumi})}
\setcounter{enumi}{2}
\tightlist
\item
  一般情形：我们可以把范数与迹的计算总结为下述命题：
\end{enumerate}

命题 4.——设 p 是 K 的特征指数，E 是 K 的一个有限次扩张。设 \(\sigma_{1}, \ldots, \sigma_{n}\) 是 E 到 K 的一个代数闭包 \(\Omega\) 中的互不相同的 K-同态，并设 \(p^{e} = [E : K]_{i}\) 。对每个 \(x \in E\) ，我们在 \(\Omega\) 中有

\[
\operatorname{Tr} _ {\mathrm{E} / \mathrm{K}} (x) = p ^ {e} \cdot \sum_ {i = 1} ^ {n} \sigma_ {i} (x), \quad \mathrm{N} _ {\mathrm{E} / \mathrm{K}} (x) = \left(\prod_ {i = 1} ^ {n} \sigma_ {i} (x)\right) ^ {p ^ {e}}.\tag{15}
\]

设 \(E_s\) 是 K 在 E 中的相对可分闭包；则 \(E_s\) 是 K 的 n 次可分扩张，且 \(\sigma_1, \ldots, \sigma_n\) 诱导 \(E_s\) 到 \(\Omega\) 中的互不相同的 K-同态；此外，E 是 \(E_s\) 的次数为 \(p^e\) 的 p-根式扩张（V，第 44 页，命题 13 及第 46 页）。于是命题 4 由公式 (6)、(7)、(13)、(14) 及 (12) 得出。

